% !Mode:: "TeX:UTF-8"
% 深圳本科毕业论文（设计）中期报告的正文，照深圳校区中期报告模板里的示范文字写。
% 节与节之间的空段照模板敲 \hitenter。模板没有参考文献一项。

\section{课题主要研究内容及进度}

\subsection{课题主要研究内容}

本课题将从以下几个方面开展研究：

\begin{enumerate}[label=\arabic*.]
  \item 人工心脏泵体数学模型的建立。……
  \item 人工心脏泵体的优化设计。……
  \item ……
\end{enumerate}

\hitenter

\subsection{进度介绍}

目前，已经完成了人工心脏泵体数学模型的建立，人工心脏泵体的优化设计。正在运用CFD软件进行相关结果的验证。

……

\section{已完成的研究工作及结果}

\subsection{人工心脏泵体的设计}

\subsubsection{叶片的参数设计}

……

结合计算的结果和结构设计的需要，调整了部分参数，基本参数如表\ref{tab:szb-1}所示。

\begin{table}[!ht]
  \centering
  \caption{叶轮主要参数表}
  \label{tab:szb-1}
  \begin{booktabs}{colspec={lll}}
    \toprule
    参数 & 数值 & 单位 \\
    \midrule
    比转速 $n_s$   & 71   & —— \\
    叶轮进口直径   & 15   & mm \\
    叶轮出口直径   & 38   & mm \\
    叶轮进口厚度   & 4    & mm \\
    叶轮出口厚度   & 4    & mm \\
    叶轮进口安装角 & 23.6 & 度 \\
    叶轮出口安装角 & 29.0 & 度 \\
    包角           & 100  & 度 \\
    \bottomrule
  \end{booktabs}
\end{table}

\hitenter

\subsubsection{叶片的三维造型}

\subsubsection{涡室的参数设计}

不加导叶的压水室大体上可分为螺旋形压水室和环形压水室两种。……

涡室的截面形状大体上有如下几种，如图\ref{fig:szb-1}所示。

\begin{figure}[!ht]
  \centering
  % 外部 TikZ 源码示例：input 在当前文档中绘图，不读取独立 PDF。
  \resizebox{14.64cm}{!}{\input{figures/volute-section}}
  \caption{涡室的截面形状}
  \label{fig:szb-1}
\end{figure}

隔舌一般指的是涡室螺旋部分的始端。第八断面与隔舌之间的夹角被称为隔舌安放角。按照经验选取，本文 $\alpha_0 = 20^\circ$。

圆形截面的计算公式如下，
\begin{equation}
  \mathit{major} = \mathit{minor} + \frac{b_3}{2} + \frac{\left| \dfrac{b_3}{2} - \mathit{minor} \right|}{2}
\end{equation}
\begin{equation}
  \mathit{height} = \mathit{minor} + \sqrt{\mathit{major}^2 - \left(\frac{b_3}{2}\right)^2}
\end{equation}
其中，$\mathit{major}$、$\mathit{minor}$ 对于椭圆段分别指的是长短轴；对于圆形段指的是半径。

……

\hitenter

\subsection{人工心脏叶轮性能优化算法}

\subsubsection{以降低血液中的切应力为目标建立优化模型}

按照血液在运动过程中，由体积相等，可得5L液体的作用时间的表达式如下，
\begin{equation}
  t = \frac{\pi b}{4Q} \left( D_2^2 - D_1^2 \right)
\end{equation}
\begin{eqdenote}
  \item $t$ & 作用时间，s；
  \item $Q$ & 流量，m$^3$/s；
  \item $b$ & 叶片的宽度，m。
\end{eqdenote}

\noindent ……

\section{后期拟完成的研究工作及进度安排}

\subsection{后期拟完成的研究工作}

后续需要完成的工作主要是……。

\hitenter

\subsection{进度安排}

具体进度安排如下：

2022年4月1日——2022年4月15日：整理已经做完的数据，分析……。

2022年4月16日——2022年4月26日：整理完成……。

2022年4月27日——2022年5月15日：……内容的研究，并完成……。

2022年5月16日——2022年6月20日：……。

\section{存在的困难及解决方案}

\subsection{存在的困难}

\begin{enumerate}[label=\arabic*.]
  \item 仿真前需要进行……。
  \item ……
  \item 有些实验设备存在购买滞后……。
\end{enumerate}

……

\hitenter

\subsection{解决方案}

针对目前存在的问题和困难，提出以下解决方案：

\begin{enumerate}[label=\arabic*.]
  \item ……
  \item ……
  \item ……
\end{enumerate}

……

\section{论文按时完成的可能性}

本课题目前已经完成了……。

目前正在进行泵体内不同工况下切应力分布图的仿真结果的统计，将要进行……。

综上所述，论文能够按期完成，并取得一定的研究成果。
